\documentclass[12pt]{article}

\usepackage[margin=1in]{geometry}
\usepackage{lmodern}
\usepackage[T1]{fontenc}
\usepackage{amsmath, amssymb, amsthm}
\newtheorem{theorem}{Theorem}
\newtheorem{definition}{Definition}
\newtheorem{proposition}{Proposition}
\newtheorem{lemma}{Lemma}
\newtheorem{corollary}{Corollary}[proposition]
\newtheorem{remark}{Remark}
\usepackage{graphicx}
\usepackage{hyperref}
\usepackage{natbib}
\usepackage{booktabs}
\usepackage{xcolor}
\usepackage{multirow}
\usepackage{longtable}
\usepackage{array}
\usepackage{microtype}
\usepackage{setspace}
\usepackage{titlesec}
\usepackage{footmisc}

%% Law-review style: section headings in small caps
\titleformat{\section}{\normalfont\large\scshape\centering}{}{0em}{}
\titleformat{\subsection}{\normalfont\normalsize\bfseries}{}{0em}{}
\titleformat{\subsubsection}{\normalfont\normalsize\itshape}{}{0em}{}

%% Spacing
\setstretch{1.5}

%% Legal case macros
\newcommand{\case}[1]{\textit{#1}}
\newcommand{\callais}{\case{Louisiana v.\ Callais}}
\newcommand{\wesberry}{\case{Wesberry v.\ Sanders}}
\newcommand{\karcher}{\case{Karcher v.\ Daggett}}
\newcommand{\shaw}{\case{Shaw v.\ Reno}}
\newcommand{\miller}{\case{Miller v.\ Johnson}}
\newcommand{\rucho}{\case{Rucho v.\ Common Cause}}
\newcommand{\lwvpa}{\case{League of Women Voters v.\ Commonwealth of Pennsylvania}}
\newcommand{\allenmilligan}{\case{Allen v.\ Milligan}}
\newcommand{\gingles}{\case{Thornburg v.\ Gingles}}
\newcommand{\harper}{\case{Harper v.\ Hall}}
\newcommand{\harkenrider}{\case{Harkenrider v.\ Hochul}}
\newcommand{\bushvvera}{\case{Bush v.\ Vera}}

%% Statute macro
\newcommand{\dia}{\textsc{Districting Integrity Act}}
\newcommand{\hh}{\textsc{Huntington-Hill}}
\newcommand{\ar}{\textsc{ApportionRegions}}
\newcommand{\usca}{\textsection}

%% Proposed statute box
\usepackage{mdframed}
\newenvironment{statute}{%
  \begin{mdframed}[linewidth=0.8pt, innerleftmargin=12pt, innerrightmargin=12pt,
                   innertopmargin=8pt, innerbottommargin=8pt,
                   backgroundcolor=gray!8]
  \small\setstretch{1.3}%
}{%
  \end{mdframed}%
}

%% Note markers
\newcommand{\est}{$^{\dagger}$}
\newcommand{\pending}{$^{\ddagger}$}

%% -----------------------------------------------------------------------

\title{%
  \textbf{From One Algorithm to One Evidence Standard}\\[0.4em]
  \large A Federal Benchmark-and-Disclosure Proposal\\[0.8em]
  \normalsize\textit{B.02 --- Advocacy Paper}%
}

\author{%
  Gio Della-Libera\\
  \textit{Apportionment Research Group}\\
  \small\texttt{giodl@microsoft.com}
}

\date{May 2026}

\begin{document}

\maketitle

\begin{abstract}
\noindent
The original version of this Article proposed that Congress mandate
\ar{} as the final-map algorithm by extending the Huntington--Hill analogy
inside States. Subsequent algorithm, legal, community, and evidence review
showed that this framing concealed rather than eliminated value judgments.
Structure, weights, resolution, seed, data corrections, VRA obligations, and
community criteria all require governance.

This revised Article proposes a narrower federal rule: every State must publish
a precommitted, partisan-input-excluded geographic benchmark and a complete
baseline-to-final decision record. The legally authorized State body or court
retains responsibility for the final plan. VRA-required changes are mandatory;
State-law, community, and error corrections are permitted only through a
publicly defined authority, alternatives, effects, and reasons.

The canonical benchmark uses census blocks, standard recursive bisection,
shared-boundary weights, and one manifest-derived seed. \ar{} remains a
research comparator, not a constitutionally compelled method. The central
federal claim is reproducibility and disclosure, not a uniquely fair output.
\end{abstract}

\noindent\textbf{Evidence and policy status.} This paper is a short internal
policy explainer, not a complete law-review constitutional analysis. It is not
the canonical national-standard
specification, enacted legislation, external constitutional review, or a
completed 50-state seed-stability certificate. The controlling candidate
decision record is
\texttt{docs/specs/2026-07-09-national-redistricting-standard-v0.1.md}.

\tableofcontents
\newpage

\section{The Revised Legislative Thesis}

\rucho{} identifies a federal constitutional justiciability limit, not a ban on
positive legislation~\citep{rucho2019}. Congress can require an auditable
congressional-election process without selecting one partisan metric or
requiring one final map.

The proposed rule has four layers:

\begin{enumerate}
\item a constitutional and data floor;
\item a reproducible geographic benchmark;
\item a governed final-plan modification process; and
\item multi-metric evaluation that does not automatically select the map.
\end{enumerate}

\section{Why The Benchmark Is Not The Final Map}

A mandatory algorithmic final map would turn one geometric preference into a
national substantive rule. It would also treat VRA, State-law, community, and
topology corrections as departures from legality rather than part of legality.

The benchmark instead creates a common evidentiary anchor. A State may depart,
but it must identify every moved block, authority, alternative, effect, public
comment, and reason. A court can review whether the published process occurred
without assuming that the benchmark itself is lawful or fair.

\section{Canonical Profile}

\begin{table}[h]
\centering
\caption{Candidate national benchmark profile.}
\begin{tabular}{ll}
\toprule
Component & Rule \\
\midrule
Unit & Census block \\
Structure & Standard floor/ceiling recursive bisection \\
Weights & Shared geographic-boundary length \\
Search & One precommitted seed \\
Seed & SHA-256 of the complete canonical input manifest \\
Population & Equality as nearly exact as practicable \\
Political/racial/COI inputs & Excluded from benchmark generation \\
\bottomrule
\end{tabular}
\end{table}

This profile differs from the original \ar{} proposal. Prime-factor structure,
county weighting, and convergence search remain testable research choices.
They are not consequences of constitutional text.

\section{Final-Plan Authority}

The final plan must change only to the extent required by the Constitution, the
Voting Rights Act, or a court order. Under \allenmilligan{}, the race-excluded
benchmark is not the appropriate vote-dilution baseline and creates no VRA
presumption, intent inference, or defense.

A final plan may also implement controlling State law, a precommitted
community-of-interest criterion, or a documented data/topology correction. The
record must include precommitted departure measures and a reproducible lower-
departure lawful alternative. Federally required remedies are compared with a
fully compliant alternative, not with proximity to the benchmark.

Where race-conscious and State-authorized partisan considerations affect
overlapping districts, the record must separate each authority, input, stage,
decision maker, and effect. Separate software runs alone do not establish
disentanglement. Partisan evidence used for racially polarized voting analysis
is not prohibited optimization input.

\section{Administration}

Congress owns assignment-affecting statutory choices. Census owns source data
and geographic corrections. NIST owns schemas, canonicalization, and technical
conformance. The Election Assistance Commission owns publication, State
assistance, public comments, and challenge intake. A balanced advisory board
includes State, technical, civil-rights, community, civic, and cross-partisan
roles but cannot amend the benchmark.

\section{Constitutional Postures}

The primary posture is an Elections Clause disclosure and process floor, with
States retaining final line-drawing authority. A federal service can publish
the benchmark when a State does not. Grants support commissions and courts. A
conditional-funding alternative remains severable if direct State duties are
invalidated.

These alternatives are not equivalent. The direct posture maximizes uniformity;
the grant and funding postures reduce commandeering risk but permit uneven
participation. The tradeoff belongs in legislative deliberation.

The principal adverse authorities are the anti-commandeering rules of
\case{Printz v.\ United States} and \case{Murphy v.\ NCAA}, Article III
standing limits under \case{TransUnion LLC v.\ Ramirez}, and the anti-coercion
limit on conditional spending under \case{NFIB v.\ Sebelius}. The model bill
therefore pairs direct duties with a federal benchmark service, a concrete-
injury standing rule, an interim-plan remedy, and a separately operative
prospective grant condition. Whether those provisions survive review is an
open legal question, not a software conclusion.

\section{Evidence Status}

BISECT now implements bounded exact split certificates, recursive certificate
trees, hostile verification corpora, tree-to-RPLAN package binding, and
pseudo-Boolean proof requests. A connected Rhode Island RCTX contains all
25,649 Census blocks and applies the established deterministic island rule.
The project has not yet generated and independently checked the three external
proofs required for a full State cut, and no national exact-readiness claim is
made. Internal review is not external peer review or legislative acceptance.

\section{Conclusion}

The revised proposal replaces ``one algorithm draws the map'' with ``one
evidence standard governs the process.'' This is less mechanically absolute
and more institutionally defensible. Its value is not removal of politics but
precommitment, reproducibility, visible alternatives, and reviewable reasons.


\bibliographystyle{plainnat}
\bibliography{references}

\end{document}
